\section{Introduction}
\label{sec:introduction}

The emergence of Time Series Foundation Models (TSFMs) has shifted forecasting paradigms from dataset-specific training toward zero-shot generalization. Models such as Chronos \cite{chronos}, Moirai \cite{moirai}, and Lag-Llama \cite{lagllama} have demonstrated remarkable ability to infer temporal patterns across diverse domains. However, standard zero-shot TSFMs struggle to synthesize information beyond their immediate context windows.

To overcome this, Retrieval-Augmented Generation (RAG)---originally popularized in natural language processing \cite{rag}---has been adapted for time series. Current state-of-the-art approaches, such as TS-RAG \cite{tsrag} and TimeRAF \cite{timeraf}, utilize deep neural encoders to project historical sequences into dense latent spaces, fusing retrieved context via learnable Adaptive Retrieval Mixers (ARM) or cross-attention mechanisms. While highly effective on dense, continuous datasets, these neural adapters break the fundamental promise of TSFMs: they require expensive pretraining on massive corpora and introduce millions of trainable parameters, rendering them computationally prohibitive for lightweight deployments.

Alternatively, parameter-free statistical retrieval (such as Euclidean k-Nearest Neighbors) offers zero-training overhead but suffers from a well-documented vulnerability: absolute scale mismatch. Because real-world time series exhibit non-stationarity and extreme magnitude variations, raw distance metrics frequently retrieve structurally dissimilar candidates, leading to catastrophic forecasting degradation.

In this work, we propose \textit{ScaleRAG}, a universal mathematical framework that enables robust retrieval augmentation without a learned projector. By explicitly normalizing historical sequences prior to indexing and mathematically restoring the retrieved continuation to the query's magnitude prior to fusion, we eliminate the need for neural encoders and adaptive mixers in the retrieval path.

We are precise about where the framework is parameter-free. Indexing, nearest-neighbour search, and scale restoration introduce \emph{zero} trainable parameters and require no pretraining: they are closed-form arithmetic over the query's own context statistics. Fusion of the restored continuation with the frozen backbone admits two variants. On dense continuous data (ETTm2) we use a fixed scalar weight, keeping the entire pipeline parameter-free. On the sparse M5 panel, where the value of retrieval varies sharply across series, we additionally employ a lightweight learned LightGBM gate over six diagnostic features (nearest-neighbour distance, retrieval disagreement, intermittency, log-volume, backbone uncertainty, and scale spread), fitted exclusively on training origins. This gate is the sole trainable component in the system and is three orders of magnitude smaller than the neural mixers it replaces; we report it explicitly rather than describing the M5 configuration as parameter-free.

Our contributions are threefold:
\begin{enumerate}
    \item We empirically isolate the \textit{Catastrophic Scale Mismatch} failure mode in non-neural retrieval, demonstrating that raw Euclidean k-NN degrades foundation model performance by a factor of 3.8$\times$ on sparse retail data.
    \item We introduce a deterministic, parameter-free scale restoration mechanism that recovers retrieval quality by +85.4\% across modalities without requiring neural retraining.
    \item We define a \textit{Taxonomy of Retrieval Regimes}, showing that neural RAG models dominate dense, continuous datasets, whereas explicit statistical restoration is the viable augmentation strategy for sparse, zero-inflated domains and resource-constrained deployments where neural adapters risk scale hallucination or cannot be trained at all. We delimit this claim with a pre-registered held-out evaluation in which our method improves squared-error accuracy over its frozen backbone yet fails to clear the pre-registered margin over strong classical baselines, and we analyse why.
\end{enumerate}
